\numoptsfalse
\yearpaper{2020}{Version A-1 \textperiodcentered{} 60 Questions \textperiodcentered{} 60 Marks \textperiodcentered{} 80 Minutes}{}
\begin{multicols}{2}

\QQ{The value of acceleration due to gravity at a height of $10\unit{km}$ from the surface of earth is $x$. At what depth inside the earth is the value of the acceleration due to gravity has the same value $x$?\ntag{D}{134}{keph107}{8}{}{Class XI Part 1 \textperiodcentered{} Ch 7}}
{\opt{$5\unit{km}$}{$20\unit{km}$}{$10\unit{km}$}{$15\unit{km}$}}

\QQ{Young's modulus of a perfect rigid body is\ntag{A}{170}{keph201}{4}{}{Class XI Part 2 \textperiodcentered{} Ch 8}\nsrc{None found in NEET/AIPMT/AIEEE/JEE papers, 2005--2026 (NCERT concept application)}}
{\opt{zero}{unity}{infinity}{between zero and unity}}

\QQ{A wheel starting from rest gains an angular velocity of $10\unit{rad/s}$ after uniformly accelerated for $5\unit{sec}$. The total angle through which it has turned is\ntag{D}{117}{keph106}{26}{}{Class XI Part 1 \textperiodcentered{} Ch 6}}
{\opt{$25\unit{rad}$}{$100\unit{rad}$}{$25\pi\unit{rad}$}{$50\pi\unit{rad}$ about a vertical axis}}

\QQ{Iceberg floats in water with part of it submerged. What is the fraction of the volume of iceberg submerged if the density of ice is $\rho_i = 0.917\unit{g\,cm^{-3}}$?\ntag{A}{181}{keph202}{2}{}{Class XI Part 2 \textperiodcentered{} Ch 9}\nsrc{None found in NEET/AIPMT/AIEEE/JEE papers, 2005--2026 (NCERT concept application)}}
{\opt{0.917}{1}{0.458}{0}}

\QQ{A sphere, a cube and a thin circular plate all of same material and same mass initially heated to same high temperature are allowed to cool down under similar conditions. Then the\ntag{D}{221}{keph203}{20}{ \textperiodcentered{} Exemplar 11.8}{Class XI Part 2 \textperiodcentered{} Ch 10}}
{\opt{plate will cool the fastest and cube the slowest.}{sphere will cool the fastest and cube the slowest.}{plate will cool the fastest and sphere the slowest.}{cube will cool the fastest and plate the slowest.}}

\QQ{In an adiabatic expansion of an ideal gas the product of pressure and volume.\ntag{A}{236}{keph204}{11}{}{Class XI Part 2 \textperiodcentered{} Ch 11}\nsrc{None found in NEET/AIPMT/AIEEE/JEE papers, 2005--2026 (NCERT concept application)}}
{\opt{Decreases}{Increases}{Remains constant}{At first increases and then decreases}}

\QQ{A certain amount of heat energy is supplied to a monoatomic ideal gas which expands at constant pressure. What fraction of the heat energy is converted into work?\ntag{A}{233}{keph204}{8}{}{Class XI Part 2 \textperiodcentered{} Ch 11}\nsrc{None found in NEET/AIPMT/AIEEE/JEE papers, 2005--2026 (NCERT concept application)}}
{\opt{1}{$\dfrac23$}{$\dfrac25$}{$\dfrac57$}}

\QQ{A tray of mass $12\unit{kg}$ is supported by two identical springs as shown in figure. When the tray is pressed down slightly and then released, it executes SHM with a time period of $1.5\unit{s}$. The spring constant of each spring is
\fig{\begin{tikzpicture}[scale=0.7]
\fill[gray!60] (-1.8,1.6) rectangle (1.8,1.72); \node[lbl,right] at (1.8,1.66) {12 kg};
\draw[decorate,decoration={coil,aspect=0.5,segment length=3pt,amplitude=4pt}] (-1.1,0.05) -- (-1.1,1.6);
\draw[decorate,decoration={coil,aspect=0.5,segment length=3pt,amplitude=4pt}] (1.1,0.05) -- (1.1,1.6);
\fill[pattern=north east lines] (-1.8,-0.12) rectangle (1.8,0); \draw (-1.8,0) -- (1.8,0);
\end{tikzpicture}}\ntag{A}{269}{keph206}{11}{}{Class XI Part 2 \textperiodcentered{} Ch 13}\nsrc{None found in NEET/AIPMT/AIEEE/JEE papers, 2005--2026 (NCERT concept application)}}
{\opt{$50\unit{N\,m^{-1}}$}{0}{$105\unit{N\,m^{-1}}$}{$\infty$}}

\QQ{A train whistling at constant frequency `$n$' is moving towards a station at a constant speed $V$. The train goes past a stationary observer on the station. The frequency `$n$' of the sound as heard by the observer is plotted as a function of time `$t$'. Identify the correct curve\ntag{X}{}{}{}{}{Doppler effect}}
{\optfig{\gaxes{$t$}{$n'$}{\draw (0.5,1.0) .. controls (0.8,1.4) and (1.2,1.4) .. (1.5,1.0);}}
{\gaxes{$t$}{$n'$}{\draw (0.2,0.2) -- (0.9,1.5) -- (1.6,0.2);}}
{\gaxes{$t$}{$n'$}{\draw (0,0.5) -- (0.9,0.5) -- (0.9,1.3) -- (1.8,1.3);}}
{\gaxes{$t$}{$n'$}{\draw (0,1.3) -- (0.9,1.3) -- (0.9,0.5) -- (1.8,0.5);}}}

\QQ{A point charge `$q$' is placed at the corner of a cube of side `$a$' as shown in the figure. What is the electric flux through the face ABCD?
\fig{\begin{tikzpicture}[scale=0.8]
\coordinate (A) at (0,1.2); \coordinate (B) at (1.2,1.2); \coordinate (G) at (1.2,0); \coordinate (H) at (0,0);
\coordinate (D) at (0.6,1.8); \coordinate (C) at (1.8,1.8); \coordinate (F) at (1.8,0.6); \coordinate (E) at (0.6,0.6);
\draw (A) -- (B) -- (G) -- (H) -- cycle; \draw (A) -- (D) -- (C) -- (B) (C) -- (F) -- (G);
\draw[densely dashed] (D) -- (E) -- (F) (E) -- (H);
\fill (A) circle (1.6pt); \node[lbl,left] at (A) {$q$};
\foreach \p/\a in {A/below right,B/below left,C/above right,D/above,E/above right,F/right,G/below,H/below} \node[font=\tiny,\a] at (\p) {\p};
\end{tikzpicture}}\ntag{A}{30}{leph101}{30}{}{Class XII Part 1 \textperiodcentered{} Ch 1}\nsrc{JEE Main 2021 (25 Feb, Shift 2) -- charge at one corner of a cube, flux through a face}}
{\opt{0}{$\dfrac{q}{24\varepsilon_0}$}{$\dfrac{q}{6\varepsilon_0}$}{$\dfrac{q}{72\varepsilon_0}$}}

\QQ{The electric field lines on the left have twice the separation on those on the right as shown in figure. If the magnitude of the field at A is $40\unit{V\,m^{-1}}$, what is the force on $20\unit{\mu C}$ charge kept at B?
\fig{\begin{tikzpicture}[scale=0.8]
\foreach \y/\yy in {0/0.15,0.6/0.45,1.2/0.75}{\draw[->] (0,\y) -- (1.5,\y) .. controls (1.9,\y) and (1.9,\yy) .. (2.3,\yy) -- (3.4,\yy);}
\fill (0.5,0.9) circle (1.4pt) node[left,lbl]{B}; \fill (2.9,0.3) circle (1.4pt) node[below,lbl]{A};
\end{tikzpicture}}\ntag{A}{20}{leph101}{20}{}{Class XII Part 1 \textperiodcentered{} Ch 1}\nsrc{None found in NEET/AIPMT/AIEEE/JEE papers, 2005--2026 (NCERT concept application)}}
{\opt{$4\times10^{-4}\unit{V\,m^{-1}}$}{$8\times10^{-4}\unit{V\,m^{-1}}$}{$16\times10^{-4}\unit{V\,m^{-1}}$}{$1\times10^{-4}\unit{V\,m^{-1}}$}}

\QQ{An infinitely long thin straight wire has uniform charge density of $\frac14\times10^{-2}\unit{C\,m^{-1}}$. What is the magnitude of electric field at a distance $20\unit{cm}$ from the axis of the wire?\ntag{D}{33}{leph101}{33}{}{Class XII Part 1 \textperiodcentered{} Ch 1}}
{\opt{$1.12\times10^8\unit{N\,C^{-1}}$}{$4.5\times10^8\unit{N\,C^{-1}}$}{$2.25\times10^8\unit{N\,C^{-1}}$}{$9\times10^8\unit{N\,C^{-1}}$}}

\QQ{A dipole of dipole moment `$P$' and moment of inertia $I$ is placed in a uniform electric field $\vec E$. If it is displaced slightly from its stable equilibrium position, the period of oscillation of dipole is\ntag{A}{27}{leph101}{27}{}{Class XII Part 1 \textperiodcentered{} Ch 1}\nsrc{None found in NEET/AIPMT/AIEEE/JEE papers, 2005--2026 (NCERT concept application)}}
{\opt{$\sqrt{\dfrac{PE}{I}}$}{$2\pi\sqrt{\dfrac{I}{PE}}$}{$\dfrac1{2\pi}\sqrt{\dfrac{PE}{I}}$}{$\pi\sqrt{\dfrac{I}{PE}}$}}

\QQ{The difference between equivalent capacitances of two identical capacitors connected in parallel to that in series is $6\unit{\mu F}$. The value of capacitance of each capacitor is\ntag{A}{72}{leph102}{28}{}{Class XII Part 1 \textperiodcentered{} Ch 2}\nsrc{None found in NEET/AIPMT/AIEEE/JEE papers, 2005--2026 (NCERT concept application)}}
{\opt{$2\unit{\mu F}$}{$3\unit{\mu F}$}{$4\unit{\mu F}$}{$6\unit{\mu F}$}}

\QQ{Figure shows three points A, B and C in a region of uniform electric field $\vec E$. The line AB is perpendicular and BC is parallel to the field lines. Then which of the following holds good? ($V_A$, $V_B$ and $V_C$ represent the electric potential at points A, B and C respectively)
\fig{\begin{tikzpicture}[scale=0.8]
\foreach \y in {0,0.5,1.0,1.5,2.0} \draw[->] (0,\y) -- (3.2,\y);
\fill (0.5,1.75) circle (1.4pt) node[above,lbl]{A}; \draw (0.5,1.75) -- (0.5,0.25);
\fill (0.5,0.25) circle (1.2pt) node[left,lbl]{B}; \draw[densely dashed] (0.5,0.25) -- (2.3,0.25); \fill (2.3,0.25) circle (1.4pt) node[right,lbl]{C};
\end{tikzpicture}}\ntag{D}{55}{leph102}{11}{}{Class XII Part 1 \textperiodcentered{} Ch 2}}
{\opt{$V_A = V_B = V_C$}{$V_A = V_B > V_C$}{$V_A = V_B < V_C$}{$V_A > V_B = V_C$}}

\QQ{When a soap bubble is charged?\ntag{A}{64}{leph102}{20}{}{Class XII Part 1 \textperiodcentered{} Ch 2}\nsrc{None found in NEET/AIPMT/AIEEE/JEE papers, 2005--2026 (NCERT concept application)}}
{\opt{Its radius increases.}{Its radius decreases.}{The radius remains the same.}{Its radius may increase or decrease.}}

\QQ{A hot filament liberates an electron with zero initial velocity. The anode potential is $1200\unit{V}$. The speed of the electron when it strikes the anode is\ntag{D}{59}{leph102}{15}{}{Class XII Part 1 \textperiodcentered{} Ch 2}}
{\opt{$1.5\times10^5\unit{m\,s^{-1}}$}{$2.5\times10^6\unit{m\,s^{-1}}$}{$2.1\times10^7\unit{m\,s^{-1}}$}{$2.5\times10^8\unit{m\,s^{-1}}$}}

\QQ{A metal rod of length $10\unit{cm}$ and a rectangular cross-section of $1\unit{cm}\times\frac12\unit{cm}$ is connected to a battery across opposite faces. The resistance will be\ntag{D}{84}{leph103}{4}{}{Class XII Part 1 \textperiodcentered{} Ch 3}}
{\opt{maximum when the battery is connected across $1\unit{cm}\times\frac12\unit{cm}$ faces.}{maximum when the battery is connected across $10\unit{cm}\times\frac12\unit{cm}$ faces.}{maximum when the battery is connected across $10\unit{cm}\times1\unit{cm}$ faces.}{same irrespective of the three faces.}}

\QQ{A car has a fresh storage battery of e.m.f.\ $12\unit{V}$ and internal resistance $2\times10^{-2}\unit{\Omega}$. If the starter motor draws a current of $80\unit{A}$. Then the terminal voltage when the starter is on is\ntag{D}{94}{leph103}{14}{}{Class XII Part 1 \textperiodcentered{} Ch 3}}
{\opt{$12\unit{V}$}{$8.4\unit{V}$}{$10.4\unit{V}$}{$9.3\unit{V}$}}

\QQ{A potentiometer has a uniform wire of length $5\unit{m}$. A battery of emf $10\unit{V}$ and negligible internal resistance is connected between its ends. A secondary cell connected to the circuit gives balancing length at $200\unit{cm}$. The emf of the secondary cell is\ntag{X}{}{}{}{}{potentiometer}}
{\opt{$4\unit{V}$}{$6\unit{V}$}{$2\unit{V}$}{$8\unit{V}$}}

\QQ{The colour code for a carbon resistor of resistance $0.28\unit{k\Omega}\pm10\%$ is\ntag{X}{}{}{}{}{resistor colour code}}
{\opt{Red, Grey, Brown, Silver}{Red, Green, Brown, Silver}{Red, Grey, Silver, Silver}{Red, Green, Silver}}

\QQ{Each resistance in the given cubical network has resistance of $1\unit{\Omega}$ and equivalent resistance between A and B is
\fig{\begin{tikzpicture}[scale=0.85]
\coordinate (a) at (0,0); \coordinate (b) at (1.8,0); \coordinate (c) at (1.8,1.8); \coordinate (d) at (0,1.8);
\coordinate (e) at (0.8,0.8); \coordinate (f) at (2.6,0.8); \coordinate (g) at (2.6,2.6); \coordinate (h) at (0.8,2.6);
\foreach \u/\v in {a/b,b/c,c/d,d/a,e/f,f/g,g/h,h/e,a/e,b/f,c/g,d/h} \draw[decorate,decoration={zigzag,segment length=3pt,amplitude=1.5pt,pre length=6pt,post length=6pt}] (\u) -- (\v);
\node[lbl,below left] at (a) {A}; \node[lbl,above right] at (g) {B};
\end{tikzpicture}}\ntag{D}{98}{leph103}{18}{ \textperiodcentered{} Example 3.5}{Class XII Part 1 \textperiodcentered{} Ch 3}}
{\opt{$\dfrac56\unit{\Omega}$}{$\dfrac65\unit{\Omega}$}{$\dfrac5{12}\unit{\Omega}$}{$\dfrac{12}5\unit{\Omega}$}}

\QQ{I-V characteristic of a copper wire of length $L$ and area of cross-section $A$ is shown in figure. The slope of the curve becomes
\fig{\gaxes{$V$}{$I$}{\draw (0.2,0.2) -- (1.6,1.5); \node[lbl,below left] at (0,0) {O};}}\ntag{A}{90}{leph103}{10}{}{Class XII Part 1 \textperiodcentered{} Ch 3}\nsrc{None found in NEET/AIPMT/AIEEE/JEE papers, 2005--2026 (NCERT concept application)}}
{\opt{More if experiment is performed at higher temperature.}{More if a wire of steel of same dimension is used.}{Less if the area of the wire is increased.}{Less if the length of the wire is increased.}}

\QQ{In the given figure, the magnetic field at `O'.
\fig{\begin{tikzpicture}[scale=0.8]
\draw (3.2,-0.75) -- (0,-0.75) arc[start angle=-90,end angle=-360,radius=0.75] -- (3.2,0);
\draw[->] (1.6,0) -- (2.4,0); \draw[->] (2.4,-0.75) -- (1.6,-0.75);
\fill (0,0) circle (1.3pt) node[above,lbl]{O}; \draw (0,0) -- (0,-0.75) node[midway,right,lbl]{$r$};
\node[lbl] at (-1.0,-0.2) {$I$};
\end{tikzpicture}}\ntag{A}{114}{leph104}{8}{}{Class XII Part 1 \textperiodcentered{} Ch 4}\nsrc{None found in NEET/AIPMT/AIEEE/JEE papers, 2005--2026 (NCERT concept application)}}
{\opt{$\dfrac34\dfrac{\mu_0I}{r}+\dfrac{\mu_0I}{4\pi r}$}{$\dfrac3{10}\dfrac{\mu_0I}{r}-\dfrac{\mu_0I}{4\pi r}$}{$\dfrac38\dfrac{\mu_0I}{r}+\dfrac{\mu_0I}{4\pi r}$}{$\dfrac38\dfrac{\mu_0I}{r}-\dfrac{\mu_0I}{4\pi r}$}}

\QQ{The magnetic field at the origin due to a current element $i\,\vec{dl}$ placed at a point with vector position $\vec r$ is\ntag{D}{113}{leph104}{7}{}{Class XII Part 1 \textperiodcentered{} Ch 4}}
{\opt{$\dfrac{\mu_0i}{4\pi}\dfrac{\vec{dl}\times\vec r}{r^3}$}{$\dfrac{\mu_0i}{4\pi}\dfrac{\vec r\times\vec{dl}}{r^3}$}{$\dfrac{\mu_0i}{4\pi}\dfrac{\vec{dl}\times\vec r}{r^2}$}{$\dfrac{\mu_0i}{4\pi}\dfrac{\vec r\times\vec{dl}}{r^2}$}}

\QQ{A long cylindrical wire of radius $R$ carries a uniform current $I$ flowing through it. The variation of magnetic field with distance `$r$' from the axis of the wire is shown by\ntag{D}{120}{leph104}{14}{ \textperiodcentered{} Example 4.7}{Class XII Part 1 \textperiodcentered{} Ch 4}}
{\optfig{\gaxes{$r$}{$B$}{\draw (0,1.2) -- (0.8,1.2) .. controls (1.0,0.6) and (1.3,0.4) .. (1.8,0.35); \draw[densely dashed] (0.8,1.2) -- (0.8,0) node[below,font=\tiny]{$r=R$};}}
{\gaxes{$r$}{$B$}{\draw[domain=0.4:1.9,samples=30] plot (\x,{0.45/\x}); \draw[densely dashed] (0.8,0.56) -- (0.8,0) node[below,font=\tiny]{$r=R$};}}
{\gaxes{$r$}{$B$}{\draw (0,0) -- (0.8,1.2) -- (1.8,0.4); \draw[densely dashed] (0.8,1.2) -- (0.8,0) node[below,font=\tiny]{$r=R$};}}
{\gaxes{$r$}{$B$}{\draw (0,0) -- (0.8,1.2) .. controls (1.0,0.6) and (1.3,0.4) .. (1.8,0.35); \draw[densely dashed] (0.8,1.2) -- (0.8,0) node[below,font=\tiny]{$r=R$};}}}

\QQ{A cyclotron is used to accelerate protons $({}^1_1\mathrm H)$, Deuterons $({}^2_1\mathrm H)$ and $\alpha$-particles $({}^4_2\mathrm{He})$. While exiting under similar conditions, the minimum K.E.\ is gained by\ntag{X}{}{}{}{}{cyclotron}}
{\opt{$\alpha$-particle}{proton}{deuteron}{same for all}}

\QQ{A paramagnetic sample shows a net magnetization of $8\unit{A\,m^{-1}}$ when placed in an external magnetic field of $0.6\unit{T}$ at a temperature of $4\unit{K}$. When the same sample is placed in an external magnetic field of $0.2\unit{T}$ at a temperature of $16\unit{K}$. The magnetization will be\ntag{X}{}{}{}{}{Curie's law}}
{\opt{$\dfrac{32}3\unit{A\,m^{-1}}$}{$\dfrac23\unit{A\,m^{-1}}$}{$6\unit{A\,m^{-1}}$}{$2.4\unit{A\,m^{-1}}$}}

\QQ{The ratio of magnetic field at the centre of a current carrying circular coil to its magnetic moment is `$x$'. If the current and the radius both are doubled. The new ratio will become\ntag{A}{127}{leph104}{21}{}{Class XII Part 1 \textperiodcentered{} Ch 4}\nsrc{None found in NEET/AIPMT/AIEEE/JEE papers, 2005--2026 (NCERT concept application)}}
{\opt{$2x$}{$4x$}{$\dfrac x4$}{$\dfrac x8$}}

\QQ{In a permanent magnet at room temperature\ntag{D}{149}{leph105}{14}{ \textperiodcentered{} Exemplar 5.3}{Class XII Part 1 \textperiodcentered{} Ch 5}}
{\opt{Magnetic moment of each molecule is zero.}{The individual molecules have non-zero magnetic moment which are all perfectly aligned.}{Domains are partially aligned.}{Domains are all perfectly aligned.}}

\QQ{A rod of length $2\unit{m}$ slides with a speed of $5\unit{m\,s^{-1}}$ on a rectangular conducting frame as shown in figure. There exists a uniform magnetic field of $0.04\unit{T}$ perpendicular to the plane of the figure. If the resistance of the rod is $3\unit{\Omega}$. The current through the rod is
\fig{\begin{tikzpicture}[scale=0.6]
\foreach \x in {0,1,...,5} \foreach \y in {0,1,2,3} \node[font=\scriptsize] at (\x,\y) {$\times$};
\draw (3.6,0.6) -- (0.5,0.6) -- (0.5,2.4) -- (3.6,2.4);
\fill[gray!70] (2.9,0.45) rectangle (3.1,2.55);
\draw[<->] (3.35,0.7) -- (3.35,2.3) node[midway,right,lbl]{2 m};
\draw[->] (2.85,1.5) -- (2.1,1.5) node[left,lbl]{$V$};
\end{tikzpicture}}\ntag{A}{163}{leph106}{10}{}{Class XII Part 1 \textperiodcentered{} Ch 6}\nsrc{None found in NEET/AIPMT/AIEEE/JEE papers, 2005--2026 (NCERT concept application)}}
{\opt{$75\unit{mA}$}{$133\unit{mA}$}{$0.75\unit{A}$}{$1.33\unit{A}$}}

\QQ{The current in a coil of inductance $0.2\unit{H}$ changes from $5\unit{A}$ to $2\unit{A}$ in $0.5\unit{sec}$. The magnitude of the average induced emf in the coil is\ntag{D}{168}{leph106}{15}{}{Class XII Part 1 \textperiodcentered{} Ch 6}}
{\opt{$0.6\unit{V}$}{$1.2\unit{V}$}{$30\unit{V}$}{$0.3\unit{V}$}}

\QQ{In the given circuit the peak voltages across C, L and R are $30\unit{V}$, $110\unit{V}$ and $60\unit{V}$ respectively. The rms value of the applied voltage is
\fig{\begin{circuitikz}[scale=0.62,transform shape,font=\Large]
\draw (0,2) to[C,l=C] (2,2) to[L,l=L] (4.2,2) to[R,l=R] (6.2,2) -- (6.2,0) to[sV] (0,0) -- (0,2);
\draw[<->] (0.2,1.2) -- (1.8,1.2) node[midway,below]{30 V}; \draw[<->] (2.1,1.2) -- (4.1,1.2) node[midway,below]{110 V}; \draw[<->] (4.3,1.2) -- (6.0,1.2) node[midway,below]{60 V};
\end{circuitikz}}\ntag{A}{187}{leph107}{11}{}{Class XII Part 1 \textperiodcentered{} Ch 7}\nsrc{None found in NEET/AIPMT/AIEEE/JEE papers, 2005--2026 (NCERT concept application)}}
{\opt{$100\unit{V}$}{$200\unit{V}$}{$70.7\unit{V}$}{$141\unit{V}$}}

\QQ{The power factor of R-L circuit is $\dfrac1{\sqrt3}$. If the inductive reactance is $2\unit{\Omega}$. The value of resistance is\ntag{A}{191}{leph107}{15}{}{Class XII Part 1 \textperiodcentered{} Ch 7}\nsrc{None found in NEET/AIPMT/AIEEE/JEE papers, 2005--2026 (NCERT concept application)}}
{\opt{$2\unit{\Omega}$}{$\sqrt2\unit{\Omega}$}{$0.5\unit{\Omega}$}{$\dfrac1{\sqrt2}\unit{\Omega}$}}

\QQ{In the given circuit, the resonant frequency is
\fig{\begin{circuitikz}[scale=0.62,transform shape,font=\Large]
\draw (0,0) to[L,l=$0.5\,\mathrm{mH}$] (0,2.2) -- (3.6,2.2) to[C,l=$20\,\mu\mathrm{F}$] (3.6,0) -- (0,0);
\end{circuitikz}}\ntag{D}{189}{leph107}{13}{}{Class XII Part 1 \textperiodcentered{} Ch 7}}
{\opt{$15.92\unit{Hz}$}{$159.2\unit{Hz}$}{$1592\unit{Hz}$}{$15910\unit{Hz}$}}

\QQ{A light beam of intensity $20\unit{W/cm^2}$ is incident normally on a perfectly reflecting surface of sides $25\unit{cm}\times15\unit{cm}$. The momentum imparted to the surface by the light per second is\ntag{X}{}{}{}{}{radiation pressure of EM waves}}
{\opt{$2\times10^{-5}\unit{kg\,m\,s^{-1}}$}{$1\times10^{-5}\unit{kg\,m\,s^{-1}}$}{$5\times10^{-5}\unit{kg\,m\,s^{-1}}$}{$1.2\times10^{-5}\unit{kg\,m\,s^{-1}}$}}

\QQ{An object approaches a convergent lens from the left of the lens with a uniform speed $5\unit{m/s}$ and stops at the focus, the image\ntag{A}{235}{leph201}{15}{ \textperiodcentered{} Exemplar 9.3}{Class XII Part 2 \textperiodcentered{} Ch 9}\nsrc{None found in NEET/AIPMT/AIEEE/JEE papers, 2005--2026 (NCERT concept application)}}
{\opt{moves away from the lens with an uniform speed $5\unit{m/s}$.}{moves away from the lens with an uniform acceleration.}{moves away from the lens with a non-uniform acceleration.}{moves towards the lens with a non-uniform acceleration.}}

\QQ{The refracting angle of a prism is A and refractive index of material of prism is $\cot\frac A2$. The angle of minimum deviation is\ntag{A}{240}{leph201}{20}{}{Class XII Part 2 \textperiodcentered{} Ch 9}\nsrc{AIPMT 2015 -- same question: prism of angle $A$ with refractive index $\cot(A/2)$, minimum deviation; also JEE Main 2022 (28 Jun, Shift 1)}}
{\opt{$180^\circ-3A$}{$180^\circ+2A$}{$90^\circ-A$}{$180^\circ-2A$}}

\QQ{The following figure shows a beam of light converging at point P. When a concave lens of focal length $16\unit{cm}$ is introduced in the path of the beam at a place shown by dotted line such that OP becomes the axis of the lens, the beam converges at a distance $x$ from the lens. The value of $x$ will be equal to
\fig{\begin{tikzpicture}[scale=0.75]
\draw (-1.5,0) -- (2.2,0); \draw[densely dashed] (0,-1.0) -- (0,1.1);
\draw[->] (-1.5,0.9) -- (-0.6,0.66); \draw (-0.6,0.66) -- (2.2,0);
\draw[->] (-1.5,-0.9) -- (-0.6,-0.66); \draw (-0.6,-0.66) -- (2.2,0);
\fill (2.2,0) circle (1.3pt) node[above,lbl]{P}; \node[lbl,below left] at (0,0) {O};
\draw[<->] (0,-0.9) -- (2.2,-0.9) node[midway,below,lbl]{12 cm};
\end{tikzpicture}}\ntag{D}{249}{leph201}{29}{ \textperiodcentered{} Exercise 9.8}{Class XII Part 2 \textperiodcentered{} Ch 9}}
{\opt{$12\unit{cm}$}{$24\unit{cm}$}{$36\unit{cm}$}{$48\unit{cm}$}}

\QQ{Three polaroid sheets $P_1$, $P_2$ and $P_3$ are kept parallel to each other such that the angle between pass axes of $P_1$ and $P_2$ is $45^\circ$ and that between $P_2$ and $P_3$ is $45^\circ$. If unpolarised beam of light of intensity $128\unit{W\,m^{-2}}$ is incident on $P_1$. What is the intensity of light coming out of $P_3$?\ntag{D}{271}{leph202}{17}{ \textperiodcentered{} Example 10.2}{Class XII Part 2 \textperiodcentered{} Ch 10}}
{\opt{$128\unit{W\,m^{-2}}$}{0}{$16\unit{W\,m^{-2}}$}{$64\unit{W\,m^{-2}}$}}

\QQ{Two poles are separated by a distance of $3.14\unit{m}$. The resolving power of human eye is 1 minute of an arc. The maximum distance from which he can identify the two poles distinctly is\ntag{X}{}{}{}{}{resolving power of the eye}}
{\opt{$10.8\unit{km}$}{$5.4\unit{km}$}{$188\unit{m}$}{$376\unit{m}$}}

\QQ{In Young's Double Slit Experiment, the distance between the slits and the screen is $1.2\unit{m}$ and the distance between the two slits is $2.4\unit{mm}$. If a thin transparent mica sheet of thickness $1\unit{\mu m}$ and R.I.\ 1.5 is introduced between one of the interfering beams, the shift in the position of central bright fringe is\ntag{A}{266}{leph202}{12}{}{Class XII Part 2 \textperiodcentered{} Ch 10}\nsrc{None found in NEET/AIPMT/AIEEE/JEE papers, 2005--2026 (NCERT concept application)}}
{\opt{$2\unit{mm}$}{$0.5\unit{mm}$}{$0.125\unit{mm}$}{$0.25\unit{mm}$}}

\QQ{The de-Broglie wavelength associated with electron of hydrogen atom in this ground state is\ntag{D}{302}{leph204}{13}{}{Class XII Part 2 \textperiodcentered{} Ch 12}}
{\opt{$0.3\unit{\text{\AA}}$}{$3.3\unit{\text{\AA}}$}{$6.26\unit{\text{\AA}}$}{$10\unit{\text{\AA}}$}}

\QQ{The following graph represents the variation of photo current with anode potential for a metal surface. Here $I_1$, $I_2$ and $I_3$ represents intensities and $\gamma_1$, $\gamma_2$, $\gamma_3$ represent frequency for curves 1, 2 and 3 respectively, then
\fig{\begin{tikzpicture}[scale=0.8]
\draw[->] (-2.4,0) -- (1.8,0) node[below left,lbl]{Anode potential};
\draw[->] (0,0) -- (0,2.0) node[above,lbl]{Photo current};
\draw (-2.0,0) .. controls (-0.8,0.1) and (-0.6,1.5) .. (0,1.5) -- (1.6,1.5); \node[font=\scriptsize] at (-1.0,1.5) {\textcircled{\tiny3}};
\draw (-1.4,0) .. controls (-0.6,0.1) and (-0.4,1.1) .. (0,1.1) -- (1.6,1.1); \node[font=\scriptsize] at (0.3,1.3) {\textcircled{\tiny2}};
\draw (-1.4,0) .. controls (-0.6,0.05) and (-0.4,0.55) .. (0,0.55) -- (1.6,0.55); \node[font=\scriptsize] at (0.3,0.3) {\textcircled{\tiny1}};
\end{tikzpicture}}\ntag{D}{279}{leph203}{6}{}{Class XII Part 2 \textperiodcentered{} Ch 11}}
{\opt{$\gamma_1 = \gamma_2$ and $I_1\ne I_2$}{$\gamma_1 = \gamma_3$ and $I_1 = I_3$}{$\gamma_1 = \gamma_2$ and $I_1 = I_2$}{$\gamma_2 = \gamma_3$ and $I_1 = I_3$}}

\QQ{The period of revolution of an electron revolving in $n^{\text{th}}$ orbit of H-atom is proportional to\ntag{A}{299}{leph204}{10}{}{Class XII Part 2 \textperiodcentered{} Ch 12}\nsrc{None found in NEET/AIPMT/AIEEE/JEE papers, 2005--2026 (NCERT concept application)}}
{\opt{$n^2$}{$\dfrac1n$}{$n^3$}{Independent of $n$}}

\QQ{Angular momentum of an electron in hydrogen atom is $\dfrac{3h}{2\pi}$ ($h$ is the Planck's constant). The K.E.\ of the electron is\ntag{D}{300}{leph204}{11}{}{Class XII Part 2 \textperiodcentered{} Ch 12}}
{\opt{$4.35\unit{eV}$}{$1.51\unit{eV}$}{$3.4\unit{eV}$}{$6.8\unit{eV}$}}

\QQ{A beam of fast moving alpha particles were directed towards a thin film of gold. The parts A, B and C of the transmitted and reflected beams corresponding to the incident parts A, B and C of the beam are shown in the adjoining diagram. The number of alpha particles in
\fig{\begin{tikzpicture}[scale=0.8]
\draw (0.5,-0.9) -- (0.5,0.6) -- (1.8,1.4) -- (1.8,-0.1) -- cycle;
\draw[->] (-1.2,0.0) -- (2.8,0.0) node[right,lbl]{A$'$}; \node[lbl,left] at (-1.2,0.0) {A};
\draw[->] (-1.2,0.25) -- (0.3,0.25); \draw (0.3,0.25) -- (0.9,0.25) arc (-90:90:0.17) -- (-1.0,0.59); \draw[->] (-0.6,0.59) -- (-1.2,0.59); \node[lbl,left] at (-1.2,0.25) {B}; \node[lbl,left] at (-1.2,0.6) {B$'$};
\draw[->] (-1.2,-0.25) -- (0.3,-0.25); \draw (0.3,-0.25) -- (0.9,-0.25) -- (2.1,-1.1); \draw[->] (1.8,-0.89) -- (2.1,-1.1) node[below right,lbl]{C$'$}; \node[lbl,left] at (-1.2,-0.25) {C};
\end{tikzpicture}}\ntag{D}{292}{leph204}{3}{}{Class XII Part 2 \textperiodcentered{} Ch 12}}
{\opt{B$'$ will be minimum and in C$'$ maximum}{A$'$ will be maximum and in C$'$ minimum}{A$'$ will be minimum and in B$'$ maximum}{C$'$ will be minimum and in B$'$ maximum}}

\QQ{Two protons are kept at a separation of $10\unit{nm}$. Let $F_n$ and $F_e$ be the nuclear force and the electromagnetic force between them\ntag{D}{313}{leph205}{8}{}{Class XII Part 2 \textperiodcentered{} Ch 13}}
{\opt{$F_e = F_n$}{$F_e\gg F_n$}{$F_e\ll F_n$}{$F_e$ and $F_n$ differ only slightly}}

\QQ{During a $\beta^-$ decay\ntag{X}{}{}{}{}{beta-decay}}
{\opt{an atomic electron is ejected.}{an electron which is already present within the nucleus is ejected.}{A neutron in the nucleus decays emitting an electron.}{A proton in the nucleus decays emitting an electron.}}

\QQ{A radio-active element has half-life of 15 years. What is the fraction that will decay in 30 years?\ntag{X}{}{}{}{}{radioactive decay law (half-life)}}
{\opt{0.25}{0.5}{0.75}{0.85}}

\QQ{A $220\unit{V}$ A.C.\ supply is connected between points A and B as shown in figure. What will be the potential difference $V$ across the capacitor?
\fig{\begin{circuitikz}[scale=0.62,transform shape,font=\Large]
\draw (0,2) node[circ]{} node[left]{A} to[empty diode] (2.5,2) -- (3.6,2);
\draw (2.5,2) to[C,l_=C] (2.5,0); \node[right] at (2.9,1.0) {$V=?$};
\draw (0,0) node[circ]{} node[left]{B} -- (3.6,0);
\node[align=center] at (0.6,1.0) {220 V\\AC};
\end{circuitikz}}\ntag{A}{339}{leph206}{17}{}{Class XII Part 2 \textperiodcentered{} Ch 14}\nsrc{None found in NEET/AIPMT/AIEEE/JEE papers, 2005--2026 (NCERT concept application)}}
{\opt{$220\unit{V}$}{$110\unit{V}$}{0}{$220\sqrt2\unit{V}$}}

\QQ{In the following circuit what are P and Q:
\fig{\begin{circuitikz}[scale=0.7,transform shape]
\draw (0,1.4) node[nand port](G1){}; \draw (0,-0.4) node[and port](G2){};
\draw (G1.in 1) -- ++(-0.6,0) node[left]{1};
\draw (G2.in 2) -- ++(-0.6,0) node[left]{0};
\draw (G1.out) -- ++(0.6,0) node[right]{P}; \draw (G2.out) -- ++(0.6,0) node[right]{Q};
\draw (G1.out) ++(0.2,0) -- ++(0,-0.4) -- ($(G2.in 1)+(-0.3,0.5)$) -- ($(G2.in 1)+(-0.3,0)$) -- (G2.in 1);
\draw (G2.out) ++(0.2,0) -- ++(0,0.4) -- ($(G1.in 2)+(-0.3,-0.5)$) -- ($(G1.in 2)+(-0.3,0)$) -- (G1.in 2);
\end{circuitikz}}\ntag{X}{}{}{}{}{digital electronics (logic gates)}}
{\opt{$P = 1$, $Q = 0$}{$P = 0$, $Q = 1$}{$P = 0$, $Q = 0$}{$P = 1$, $Q = 1$}}

\QQ{A positive hole in a semiconductor is\ntag{D}{327}{leph206}{5}{}{Class XII Part 2 \textperiodcentered{} Ch 14}}
{\opt{an anti-particle of electron.}{a vacancy created when an electron leaves a covalent bond.}{absence of free electrons.}{an artificially created particle.}}

\QQ{Two long straight parallel wires are a distance $2d$ apart. They carry steady equal currents flowing out of the plane of the paper. The variation of magnetic field B along the line $xx'$ is given by\ntag{A}{123}{leph104}{17}{}{Class XII Part 1 \textperiodcentered{} Ch 4}\nsrc{None found in NEET/AIPMT/AIEEE/JEE papers, 2005--2026 (NCERT concept application)}}
{\optfig{\bwires{1}{1}{-1}{-1}}{\bwires{-1}{-1}{1}{1}}{\bwires{-1}{1}{1}{-1}}{\bwires{-1}{-1}{-1}{-1}}}

\QQ{A cylindrical wire has a mass $(0.3\pm0.003)\unit{g}$, radius $(0.5\pm0.005)\unit{mm}$ and length $(6\pm0.06)\unit{cm}$. The maximum percentage error in the measurement of its density is\ntag{A}{6}{keph101}{6}{}{Class XI Part 1 \textperiodcentered{} Ch 1}\nsrc{None found in NEET/AIPMT/AIEEE/JEE papers, 2005--2026 (NCERT concept application)}}
{\opt{1}{2}{3}{4}}

\QQ{At a metro station, a girl walks up a stationary escalator in $20\unit{sec}$. If she remains stationary on the escalator, then the escalator take her up in $30\unit{sec}$. The time taken by her to walk up on the moving escalator will be\ntag{A}{46}{keph103}{20}{}{Class XI Part 1 \textperiodcentered{} Ch 3}\nsrc{NEET 2017 -- same question: girl walking up a stationary and a moving escalator}}
{\opt{$25\unit{sec}$}{$60\unit{sec}$}{$12\unit{sec}$}{$10\unit{sec}$}}

\QQ{Rain is falling vertically with a speed of $12\unit{m\,s^{-1}}$. A woman rides a bicycles with a speed of $12\unit{m\,s^{-1}}$ in east to west direction. What is the direction in which she should hold her umbrella?\ntag{D}{31}{keph103}{5}{ \textperiodcentered{} Example 3.1}{Class XI Part 1 \textperiodcentered{} Ch 3}}
{\opt{$30^\circ$ towards East}{$45^\circ$ towards East}{$30^\circ$ towards West}{$45^\circ$ towards West}}

\QQ{One end of a string of length `$l$' is connected to a particle of mass `$m$' and the other to a small peg on a smooth horizontal table. If the particle moves in a circle with speed `$v$', the net force on the particle (directed towards the centre) is: ($T$ is the tension in the string)\ntag{D}{69}{keph104}{21}{ \textperiodcentered{} Exercise 4.4}{Class XI Part 1 \textperiodcentered{} Ch 4}}
{\opt{$T$}{$T-\dfrac{mv^2}{l}$}{$T+\dfrac{mv^2}{l}$}{0}}

\QQ{A body is initially at rest. It undergoes one-dimensional motion with constant acceleration. The power delivered to it at time `$t$' is proportional to\ntag{D}{90}{keph105}{20}{ \textperiodcentered{} Exercise 5.9}{Class XI Part 1 \textperiodcentered{} Ch 5}}
{\opt{$t^{1/2}$}{$t$}{$t^{3/2}$}{$t^2$}}

\QQ{A thin uniform rectangular plate of mass $2\unit{kg}$ is placed in X-Y plane as shown in figure. The moment of inertia about $x$-axis is $I_x = 0.2\unit{kg\,m^2}$ and the moment of inertia about $y$-axis is $I_y = 0.3\unit{kg\,m^2}$. The radius of gyration of the plate about the axis passing through O and perpendicular to the plane of the plate is
\fig{\begin{tikzpicture}[scale=0.75]
\draw (-1.3,-0.5) -- (0.6,-0.5) -- (1.3,0.5) -- (-0.6,0.5) -- cycle;
\draw[->] (-1.6,0) -- (1.9,0) node[right,lbl]{X};
\draw[->] (0,-0.2) -- (0,1.3) node[right,lbl]{Z};
\draw[->] (0.9,1.2) -- (-0.9,-1.2) node[left,lbl]{Y};
\node[lbl,above left] at (0,0) {O};
\end{tikzpicture}}\ntag{X}{}{}{}{}{perpendicular axis theorem}}
{\opt{$50\unit{cm}$}{$5\unit{cm}$}{$38.7\unit{cm}$}{$31.6\unit{cm}$}}

%%BEGIN-ANALYSIS
\end{multicols}
\clearpage
\begin{center}{\pgfornament[width=5cm,color=gilt]{87}}\par
{\fontsize{20}{24}\selectfont\color{navy}\scshape KCET 2020 \textperiodcentered{} Paper Analysis}\par
{\color{gilt}\pgfornament[width=5cm,color=gilt,symmetry=h]{87}}\end{center}
\vspace{4pt}
\begin{multicols}{2}
\textbf{\color{crimson}1. Direct vs Applied}\par\smallskip
{\footnotesize Checked question by question against the four 2026-27 NCERT books (XI Part 1 \& 2, XII Part 1 \& 2). Out of 60, 49 questions are in the present syllabus; percentages are of these.}\par\smallskip
\begin{tabular}{@{}lrr@{}}\toprule
Direct from NCERT & \textbf{25} & 51\%\\
\quad text, summary, table or figure & 18 &\\
\quad NCERT solved Example & 4 &\\
\quad NCERT Exercise & 3 &\\
Applied (concept in NCERT) & \textbf{24} & 49\%\\\midrule
Analysed & 49 &\\
Deleted portion (not analysed) & 11 &\\\midrule
\textbf{Total} & \textbf{60} &\\\bottomrule\end{tabular}\par\smallskip
{\footnotesize Direct: Q 1, 3, 5, 12, 15, 17, 18, 19, 22, 25, 26, 30, 32, 35, 39, 40, 43, 44, 46, 47, 48, 53, 57, 58, 59.\\ From NCERT Examples: Q 22 (Example 3.5), 26 (Example 4.7), 40 (Example 10.2), 57 (Example 3.1).\\ From NCERT Exercises: Q 39 (Exercise 9.8), 58 (Exercise 4.4), 59 (Exercise 5.9).}\par\medskip
\textbf{\color{crimson}2. Where the applied questions resemble earlier exams}\par\smallskip
Of the 24 applied questions, \textbf{3} closely resemble a question from an earlier national exam:
NEET/AIPMT/AIIMS 2, JEE Main/Advanced 1, AIEEE 0. The other 21 (Q 2, 4, 6, 7, 8, 11, 13, 14, 16, 23, 24, 29, 31, 33, 34, 37, 42, 45, 51, 54, 55) had no close match in the NEET, AIPMT, AIIMS, AIEEE and JEE papers of 2005--2026 that were searched; they are straightforward applications of NCERT concepts.\par
{\footnotesize Strongest repeats: Q38 (AIPMT 2015, identical), Q56 (NEET 2017, identical), Q10 (JEE Main 2021).}\par\medskip
\textbf{\color{crimson}3. Difficulty level}\par\smallskip
\begin{tabular}{@{}lr@{}}\toprule
Easy & 29\\ Moderate & 19\\ Hard & 1\\\bottomrule\end{tabular}\par\smallskip
{\footnotesize Hard: Q 24. Overall: \textbf{easy to moderate}; Q24 (field at O due to straight and curved parts) is the only lengthy calculation. Class XI: 13 questions, Class XII: 36 (analysed questions only). (Difficulty is the editor's assessment.)}\par\medskip
\columnbreak
\textbf{\color{crimson}4. Chapter-wise distribution}\par\smallskip
{\footnotesize\begin{tabular}{@{}p{3.6cm}c@{\,}r@{\,}r@{\,}r@{}}\toprule
Chapter & Cl. & Qs & Dir & App\\\midrule
Units \& Measurement & XI & 1 & 0 & 1\\
Motion in a Plane & XI & 2 & 1 & 1\\
Laws of Motion & XI & 1 & 1 & 0\\
Work, Energy \& Power & XI & 1 & 1 & 0\\
Rotational Motion & XI & 1 & 1 & 0\\
Gravitation & XI & 1 & 1 & 0\\
Mech.\ Prop.\ of Solids & XI & 1 & 0 & 1\\
Mech.\ Prop.\ of Fluids & XI & 1 & 0 & 1\\
Thermal Prop.\ of Matter & XI & 1 & 1 & 0\\
Thermodynamics & XI & 2 & 0 & 2\\
Oscillations & XI & 1 & 0 & 1\\
Electric Charges \& Fields & XII & 4 & 1 & 3\\
Electrostatic Potential \& Capacitance & XII & 4 & 2 & 2\\
Current Electricity & XII & 4 & 3 & 1\\
Moving Charges \& Magnetism & XII & 5 & 2 & 3\\
Magnetism \& Matter & XII & 1 & 1 & 0\\
Electromagnetic Induction & XII & 2 & 1 & 1\\
Alternating Current & XII & 3 & 1 & 2\\
Ray Optics & XII & 3 & 1 & 2\\
Wave Optics & XII & 2 & 1 & 1\\
Dual Nature & XII & 1 & 1 & 0\\
Atoms & XII & 4 & 3 & 1\\
Nuclei & XII & 1 & 1 & 0\\
Semiconductor Electronics & XII & 2 & 1 & 1\\
\bottomrule\end{tabular}}\par\medskip
\textbf{\color{crimson}5. Teaching takeaways}\par\smallskip
{\footnotesize\begin{itemize}[leftmargin=1em,nosep]
\item Solved Examples asked again: Example 3.5 (cube of twelve 1\,$\Omega$ resistors, $R=5/6\,\Omega$), Example 3.1 (rain and umbrella), Example 4.7 (field inside and outside a thick wire), Example 10.2 (three polaroids).
\item Exercises repeated: Exercise 4.4 (particle on a string tied to a peg), Exercise 5.9 ($P\propto t$ at constant acceleration), Exercise 9.8 (concave lens in a converging beam).
\item Eleven questions are on deleted portions (Doppler effect, potentiometer, colour code, cyclotron, Curie's law, radiation pressure, resolving power, $\beta$-decay, half-life, logic gates, perpendicular-axis theorem).
\end{itemize}}\par\medskip
\textbf{\color{crimson}6. Deleted portions}\par\smallskip
\textbf{11} of the 60 questions are on topics removed from the 2026-27 NCERT (rationalised syllabus). They are marked \textsc{Deleted Portion} and are not analysed: Doppler effect (Q 9); potentiometer (Q 20); resistor colour code (Q 21); cyclotron (Q 27); Curie's law (Q 28); radiation pressure of EM waves (Q 36); resolving power of the eye (Q 41); beta-decay (Q 49); radioactive decay law (half-life) (Q 50); digital electronics (logic gates) (Q 52); perpendicular axis theorem (Q 60).\par
\end{multicols}
%%END-ANALYSIS



